\documentclass[10pt,a4paper]{amsart}
\usepackage{amsmath,amssymb,amsthm,graphicx}
\usepackage{verbatim}
\newenvironment{reference}{\comment}{\endcomment}
%\newenvironment{reference}{}{}
\newenvironment{slogan}{\comment}{\endcomment}
\newenvironment{history}{\comment}{\endcomment}

% For commutative diagrams we use Xy-pic
\usepackage[all]{xy}

% We use 2cell for 2-commutative diagrams.
\xyoption{2cell}
% Theorem environments.
%
\theoremstyle{plain}
\newtheorem{theorem}[subsection]{Theorem}
\newtheorem{proposition}[subsection]{Proposition}
\newtheorem{lemma}[subsection]{Lemma}

\theoremstyle{definition}
\newtheorem{definition}[subsection]{Definition}
\newtheorem{example}[subsection]{Example}
\newtheorem{exercise}[subsection]{Exercise}
\newtheorem{situation}[subsection]{Situation}

\theoremstyle{remark}
\newtheorem{remark}[subsection]{Remark}
\newtheorem{remarks}[subsection]{Remarks}

\numberwithin{equation}{subsection}

% Macros
%
\def\lim{\mathop{\mathrm{lim}}\nolimits}
\def\colim{\mathop{\mathrm{colim}}\nolimits}
\def\Spec{\mathop{\mathrm{Spec}}}
\def\Hom{\mathop{\mathrm{Hom}}\nolimits}
\def\Ext{\mathop{\mathrm{Ext}}\nolimits}
\def\SheafHom{\mathop{\mathcal{H}\!\mathit{om}}\nolimits}
\def\SheafExt{\mathop{\mathcal{E}\!\mathit{xt}}\nolimits}
\def\Sch{\mathit{Sch}}
\def\Mor{\mathop{\mathrm{Mor}}\nolimits}
\def\Ob{\mathop{\mathrm{Ob}}\nolimits}
\def\Sh{\mathop{\mathit{Sh}}\nolimits}
\def\NL{\mathop{N\!L}\nolimits}
\def\CH{\mathop{\mathrm{CH}}\nolimits}
\def\proetale{{pro\text{-}\acute{e}tale}}
\def\etale{{\acute{e}tale}}
\def\QCoh{\mathit{QCoh}}
\def\Ker{\mathop{\mathrm{Ker}}}
\def\Im{\mathop{\mathrm{Im}}}
\def\Coker{\mathop{\mathrm{Coker}}}
\def\Coim{\mathop{\mathrm{Coim}}}

% Boxtimes
%
\DeclareMathSymbol{\boxtimes}{\mathbin}{AMSa}{"02}

%
% Macros for moduli stacks/spaces
%
\def\QCohstack{\mathcal{QC}\!\mathit{oh}}
\def\Cohstack{\mathcal{C}\!\mathit{oh}}
\def\Spacesstack{\mathcal{S}\!\mathit{paces}}
\def\Quotfunctor{\mathrm{Quot}}
\def\Hilbfunctor{\mathrm{Hilb}}
\def\Curvesstack{\mathcal{C}\!\mathit{urves}}
\def\Polarizedstack{\mathcal{P}\!\mathit{olarized}}
\def\Complexesstack{\mathcal{C}\!\mathit{omplexes}}
% \Pic is the operator that assigns to X its picard group, usage \Pic(X)
% \Picardstack_{X/B} denotes the Picard stack of X over B
% \Picardfunctor_{X/B} denotes the Picard functor of X over B
\def\Pic{\mathop{\mathrm{Pic}}\nolimits}
\def\Picardstack{\mathcal{P}\!\mathit{ic}}
\def\Picardfunctor{\mathrm{Pic}}
\def\Deformationcategory{\mathcal{D}\!\mathit{ef}}
\usepackage{cleveref}
\crefname{originallemma}{original lemma}{original lemmas}
\Crefname{originallemma}{Original lemma}{Original lemmas}
\makeatletter
\newlabel{derived-lemma-commutes-with-countable-sums}{{094A}{source}{}{sourcepointer-1}{}}
\newlabel{derived-lemma-commutes-with-countable-sums@cref}{{[originallemma][0][]094A}{[1][0][]source}{}{}{}}
\makeatother
\crefname{originallemma}{original lemma}{original lemmas}
\Crefname{originallemma}{Original lemma}{Original lemmas}
\makeatletter
\newlabel{derived-lemma-factor-through}{{09SP}{source}{}{sourcepointer-2}{}}
\newlabel{derived-lemma-factor-through@cref}{{[originallemma][0][]09SP}{[1][0][]source}{}{}{}}
\makeatother
\crefname{originallemma}{original lemma}{original lemmas}
\Crefname{originallemma}{Original lemma}{Original lemmas}
\makeatletter
\newlabel{derived-lemma-write-as-colimit}{{09SN}{source}{}{sourcepointer-3}{}}
\newlabel{derived-lemma-write-as-colimit@cref}{{[originallemma][0][]09SN}{[1][0][]source}{}{}{}}
\makeatother
\crefname{originalreference}{original reference}{original references}
\Crefname{originalreference}{Original reference}{Original references}
\makeatletter
\newlabel{lemma-generator-with-support}{{0A9A}{source}{}{sourcepointer-4}{}}
\newlabel{lemma-generator-with-support@cref}{{[originalreference][0][]0A9A}{[1][0][]source}{}{}{}}
\makeatother
\crefname{originalreference}{original reference}{original references}
\Crefname{originalreference}{Original reference}{Original references}
\makeatletter
\newlabel{lemma-quasi-coherence-internal-hom}{{0A6H}{source}{}{sourcepointer-5}{}}
\newlabel{lemma-quasi-coherence-internal-hom@cref}{{[originalreference][0][]0A6H}{[1][0][]source}{}{}{}}
\makeatother
\title{Maps from pseudo-coherent objects into supported complexes kill a perfect function supported off T}
\author{The Stacks Project Authors}
\date{stacks-project:spaces-perfect.tex:local-snapshot-20261012; GNU FDL 1.2 or later}
\begin{document}
\maketitle
\noindent\textit{Editorial note.} All delivered mathematical text is the author's own, copied from the pinned Stacks Project chapter file spaces-perfect.tex; no statement or proof step was rewritten, reordered or removed.\par\medskip
\noindent\textit{Editorial note.} The counted proof is the author's complete argument: the proof is complete: the reduction to O\_X[0], the finite-stage realization and the cone construction\par\medskip
\noindent\textit{Editorial note.} The label lemma-map-from-pseudo-coherent-to-complex-with-support names the counted statement, and the proof environment immediately following it in the source is delivered verbatim as its proof.\par\medskip
\noindent\textit{Editorial note.} The counted statement is carried by the author's own lemma environment in the chapter source.\par\medskip
\noindent\textit{Allowed standard prerequisites (editorial).}\begin{itemize}
\item the local subcategory of D\_QCoh supported on a closed subset and its perfect generator
\item internal Hom in D\_QCoh, its quasi-coherence and compatibility with restriction to opens
\item colimit representations using the generator and compactness of O\_X
\item the factor-through lemma for maps into the thick subcategory generated by E
\item perfectness of the cone over a map to a perfect object supported on T
\end{itemize}\medskip
\ifdefined\hypertarget\hypertarget{sourcepointer-1}{}\fi
\noindent\textit{Source reference.} Original lemma 094A: sources/stacks-chow/tags.txt, source lines 11543--11543. Prerequisite pointer invoked at proof line(s) 4258; it names a standard fact of a sibling chapter that the counted argument uses, and every step of the counted proof built on it is retained verbatim in the delivered body.\par\medskip
\ifdefined\hypertarget\hypertarget{sourcepointer-2}{}\fi
\noindent\textit{Source reference.} Original lemma 09SP: sources/stacks-chow/tags.txt, source lines 12362--12362. Prerequisite pointer invoked at proof line(s) 4259; it names a standard fact of a sibling chapter that the counted argument uses, and every step of the counted proof built on it is retained verbatim in the delivered body.\par\medskip
\ifdefined\hypertarget\hypertarget{sourcepointer-3}{}\fi
\noindent\textit{Source reference.} Original lemma 09SN: sources/stacks-chow/tags.txt, source lines 12361--12361. Prerequisite pointer invoked at proof line(s) 4252; it names a standard fact of a sibling chapter that the counted argument uses, and every step of the counted proof built on it is retained verbatim in the delivered body.\par\medskip
\ifdefined\hypertarget\hypertarget{sourcepointer-4}{}\fi
\noindent\textit{Source reference.} Original reference 0A9A: sources/stacks-spaces-perfect/pinned-original-source.tex, source lines 3997--3997. Prerequisite pointer invoked at proof line(s) 4248; it names a standard fact of the same chapter that the counted argument uses, and every step of the counted proof built on it is retained verbatim in the delivered body.\par\medskip
\ifdefined\hypertarget\hypertarget{sourcepointer-5}{}\fi
\noindent\textit{Source reference.} Original reference 0A6H: sources/stacks-spaces-perfect/pinned-original-source.tex, source lines 3499--3499. Prerequisite pointer invoked at proof line(s) 4238; it names a standard fact of the same chapter that the counted argument uses, and every step of the counted proof built on it is retained verbatim in the delivered body.\par\medskip
\begin{lemma}
\label{lemma-map-from-pseudo-coherent-to-complex-with-support}
Let $S$ be a scheme. Let $X$ be a quasi-compact and quasi-separated
algebraic space over $S$. Let $T \subset |X|$
be a closed subset such that the complement $U \subset X$ is quasi-compact.
Let $\alpha : P \to E$ be a morphism of $D_\QCoh(\mathcal{O}_X)$ with
either
\begin{enumerate}
\item $P$ is perfect and $E$ supported on $T$, or
\item $P$ pseudo-coherent, $E$ supported on $T$, and $E$ bounded below.
\end{enumerate}
Then there exists a perfect complex of $\mathcal{O}_X$-modules $I$
and a map $I \to \mathcal{O}_X[0]$ such that
$I \otimes^\mathbf{L} P \to E$ is zero and such that
$I|_U \to \mathcal{O}_U[0]$ is an
isomorphism.
\end{lemma}

\begin{proof}
Set $\mathcal{D} = D_{\QCoh, T}(\mathcal{O}_X)$. In both cases the complex
$K = R\SheafHom(P, E)$ is an object of $\mathcal{D}$. See
Lemma \ref{lemma-quasi-coherence-internal-hom} for quasi-coherence.
It is clear that $K$ is supported on $T$ as formation of $R\SheafHom$
commutes with restriction to opens.
The map $\alpha$ defines an element of
$H^0(K) = \Hom_{D(\mathcal{O}_X)}(\mathcal{O}_X[0], K)$.
Then it suffices to prove the result for the map
$\alpha : \mathcal{O}_X[0] \to K$.

\medskip\noindent
Let $E \in \mathcal{D}$ be a perfect generator, see
Lemma \ref{lemma-generator-with-support}. Write
$$
K = \text{hocolim} K_n
$$
as in Derived Categories, Lemma \ref{derived-lemma-write-as-colimit}
using the generator $E$. Since the functor $\mathcal{D} \to D(\mathcal{O}_X)$
commutes with direct sums, we see that $K = \text{hocolim} K_n$
holds in $D(\mathcal{O}_X)$. Since $\mathcal{O}_X$ is a compact
object of $D(\mathcal{O}_X)$ we find an $n$ and a morphism
$\alpha_n : \mathcal{O}_X \to K_n$ which gives rise to $\alpha$, see
Derived Categories, Lemma \ref{derived-lemma-commutes-with-countable-sums}.
By Derived Categories, Lemma \ref{derived-lemma-factor-through}
applied to the morphism $\mathcal{O}_X[0] \to K_n$ in the ambient
category $D(\mathcal{O}_X)$ we see that $\alpha_n$ factors as
$\mathcal{O}_X[0] \to Q \to K_n$ where $Q$ is an object
of $\langle E \rangle$. We conclude that $Q$ is a perfect complex
supported on $T$.

\medskip\noindent
Choose a distinguished triangle
$$
I \to \mathcal{O}_X[0] \to Q \to I[1]
$$
By construction $I$ is perfect, the map $I \to \mathcal{O}_X[0]$
restricts to an isomorphism over $U$, and the composition
$I \to K$ is zero as $\alpha$ factors through $Q$.
This proves the lemma.
\end{proof}

\end{document}
